\section{Related Work}\label{sec:related}
%==============================================================================

\subsection{Zero-Error and Side-Information Lineage}\label{sec:related-source-coding}

Shannon's zero-error framework and its graph-theoretic refinements by Körner and Lovász provide the closest conceptual starting point~\cite{shannon1956zero,korner1973graphs,lovasz1979shannon}. Körner--Orlitsky's zero-error synthesis is also a direct reference point for this lineage~\cite{korner2002zero}. Witsenhausen's zero-error side-information problem is directly relevant because it studies exact recovery under side information through graph coloring and label budgets~\cite{witsenhausen1976zero}. Our one-shot coloring statements lie in that lineage. As in Witsenhausen's formulation, an observation law induces the relevant confusability graph. The difference is that here the observation law is generated by a deterministic multi-fact tuple-space architecture with admissible coordinate-subset views, so the model itself constrains what graphs can arise.

Slepian--Wolf coding and classical entropy converses supply the second line of influence~\cite{slepian1973noiseless,fano1961transmission,witsenhausen1975conditional,cover2006elements}. Coding-for-computing, functional compression, and graph-entropy/characteristic-graph viewpoints are also close neighbors because they study deterministic decoder targets and coding questions once an underlying graph or function structure has been specified~\cite{orlitsky2001coding,alon2002source,doshi2010functional,korner1973graphs,simonyi2001graph}. Recent zero-error function-compression work also sits in this neighborhood~\cite{guang2023zero,liu2021leakage}.

Witsenhausen's setting starts from a side-information law and studies the coding problem for the graph it induces. Our setting starts from a deterministic partial-view architecture on latent tuples, with observations obtained by admissible coordinate-subset projections, and asks what graph structure that architecture can generate and what recovery laws follow from it. In the exact full-tuple-space coordinate-view model analyzed here, the realizable confusability relations are precisely those determined by upward-closed families of coordinate-agreement sets. Characteristic-graph formalisms handle arbitrary deterministic functions of the source, whereas the coordinate-projection restriction used here also yields coordinatewise alphabet-permutation automorphisms, a monotone agreement-set graph class, and closure under block composition by strong product. The paper therefore does not claim a new theorem about arbitrary graphs; it identifies a deterministic partial-view model with a fully characterized model-generated graph class on the full labeled tuple space. What it does not attempt is the broader unlabeled-isomorphism classification, or the corresponding classifications for restricted realized-state families and stochastic support-overlap models.

In our setting the side information is deterministic rather than stochastic, but it plays the same operational role: exact recovery becomes impossible when the available observation/tag alphabet is too small to separate the remaining alternatives.

After a confusability graph is fixed, the colorability, strong-power, Shannon-capacity, and Lov\'asz-$\vartheta$ machinery used later is classical. The novelty here is identifying transitivity of view-generated confusability as the model-side condition that produces the cluster-graph equality subclass, together with meet-witnessing and fiber coherence as structural sufficient conditions. On the upper-bound side, this follows the same classical-to-modern arc that includes Haemers-type bounds and asymptotic-spectrum viewpoints~\cite{haemers1978upper,zuiddam2019asymptotic,deboer2024asymptotic}.

\subsection{Consistency-Constrained Storage and Interaction}\label{sec:related-distributed}

Consistency-constrained storage problems, including multi-version coding and related distributed-storage converses, show that correctness requirements impose unavoidable information costs~\cite{rashmi2015multiversion}. The consistency literature establishes fundamental tradeoffs between consistency, availability, and partition tolerance~\cite{brewer2000cap,gilbert2002cap,gray1978notes}, while correctness models such as linearizability specify formal guarantees for concurrent operations~\cite{herlihy1990linearizability}. Here the object under study is a modifiable encoding architecture, and the main question is the exact-consistency cost of multiple independently writable representations of one latent fact.

\subsection{Computational Realizations}\label{sec:related-meta}

Reflection, metadata systems, and maintenance mechanisms in host platforms provide examples of how the boundary-case realizability conditions can be instantiated in practice~\cite{kiczales1991art,smith1984reflection}. These examples are secondary to the paper's main zero-error and graph-capacity contribution.

\subsection{Formal Verification}\label{sec:related-formal}

The Lean 4 formalization places the work in the tradition of mechanized mathematics and verified theory development~\cite{demoura2021lean4,mathlib2020,leroy2009compcert}.

%==============================================================================
